\section{Toward a finite basis}\label{app:basis}

A finite basis for full owned choice or its intrinsic core is open. The polynomial guard algorithm of \cref{sec:one-owner} does not supply one for the four-element expansion, whose abstract guard tables also lack a realizability criterion analogous to reset's: canonical comparison of tables extracted from terms does not characterize which tables can occur. Since
that generator is conservative, it may yield to the methods that classify the conservative models of reset. The polynomial bounds of \cref{sec:one-owner} concern trees, and shared subexpressions require a separate size analysis.

There is a conditional route to a finite basis. Each owned choice is a semilattice operation, and for fixed $K$ the full finite-signature variety is congruence meet-semidistributive. Concretely, the semilattice witness terms displayed just before Willard's Theorem~2.1 satisfy his condition $C_2(T)$ when their product is interpreted as any fixed $\star_r$
\cite{Willard2000}; the characterization then supplies meet-semidistributivity, and the expansion by other operations preserves those term identities. Willard's finite-basis theorem applies if the variety has a finite residual bound \cite{Willard2000}. Classifying quotients of the generator does not bound irreducibles throughout the variety, and reset's residual
largeness does not transfer automatically to an expansion by other operations.

Two specific obstacles stand in the way of a residual bound. The pair test domain is not a subalgebra, and testing all identities on that domain does not bound the irreducible models of those identities. Likewise, the $K+1$-state player quotients describe the generator's reduct, not every subalgebra and quotient of arbitrary powers. In the opposite direction, the reset
combs would prove residual largeness for full owned choice if each could be expanded to a member of the full generated variety with the prescribed reset term operation. Such an expansion would remain subdirectly irreducible: every nonidentity congruence of the expansion is a reset congruence and therefore contains the comb's nontrivial monolith. The missing step is
constructing expansions satisfying all full-signature identities.

The tractable fragments are also constrained by symmetry. With role switching, one owned choice defines every other:
\[
 x\star_b y=\pi(\pi^{-1}x\star_a\pi^{-1}y),\qquad\pi(a)=b.
\]
So a role-symmetric fragment containing unrestricted owned choice already contains the hard core, and a tractable role-symmetric fragment needs something weaker.
